\documentclass[11pt,a4paper]{article}

% ---------- Paquetes ----------
\usepackage[utf8]{inputenc}
\usepackage[spanish]{babel}
\usepackage[T1]{fontenc}
\usepackage{amsmath}
\usepackage{graphicx}
\usepackage{tikz}
\usetikzlibrary{arrows.meta}
\usepackage[margin=2.5cm]{geometry}
\usepackage[style=ieee]{biblatex}
\addbibresource{referencias.bib}

\title{PE-U3: Fragmentación y Tolerancia a Fallos en Bases de Datos Distribuidas \\ \large Sistema de Gestión Académica (SGA) sobre CockroachDB}
\author{Equipo BCEL}
\date{\today}

\begin{document}
\maketitle

% ============================================================
% NOTA: Aquí el equipo va agregando las demás secciones
% (introducción, fragmentación, benchmarks, etc.)
% ============================================================

\section{Prueba de tolerancia a fallos}

\subsection{Control de concurrencia distribuido y consenso}

En una base de datos distribuida, múltiples transacciones concurrentes pueden acceder
a los mismos datos replicados en distintos nodos, lo que puede producir anomalías como
lecturas sucias o actualizaciones perdidas si no existe un mecanismo de control
\cite{ozsu2020}. El protocolo de bloqueo en dos fases (2PL) resuelve este problema
exigiendo que toda transacción adquiera sus bloqueos en una fase de crecimiento y los
libere en una fase de reducción, sin volver a adquirir ninguno. Su variante estricta
(S2PL) mantiene los bloqueos exclusivos hasta el \emph{commit}, garantizando
seriabilidad y recuperabilidad. En el entorno distribuido, los bloqueos se coordinan
entre los nodos que albergan cada fragmento.

El bloqueo introduce el riesgo de interbloqueo (deadlock). Su detección distribuida se
realiza construyendo grafos de espera (wait-for graphs) locales que luego se combinan;
un ciclo en el grafo global indica un deadlock, que se resuelve abortando una de las
transacciones implicadas \cite{ozsu2020}.

Para garantizar la atomicidad de una transacción que abarca varios nodos se emplea el
protocolo de commit en dos fases (2PC): un coordinador envía una fase de preparación y,
si todos los participantes responden afirmativamente, ordena el commit definitivo. Su
debilidad es el bloqueo si el coordinador falla tras la preparación; el protocolo 3PC
añade una fase intermedia para reducir ese bloqueo, a costa de más mensajes
\cite{vansteen2023,skeen1981}.

Los sistemas modernos como CockroachDB sustituyen la fragilidad del 2PC clásico por el
algoritmo de consenso \textbf{Raft} \cite{ongaro2014}. Raft elige un líder por cada
grupo de réplicas (rango); el líder recibe las escrituras, las anexa a un log y las
replica a los seguidores. Una entrada se considera confirmada (committed) cuando la
mayoría de las réplicas la ha almacenado, es decir, cuando se alcanza el quórum. Si el
líder cae, los seguidores inician una elección y eligen uno nuevo, garantizando
continuidad y seguridad: nunca se pierde una entrada confirmada \cite{taft2020}.

\subsection{Diagrama de secuencia: Raft ante la caída de un nodo}

\begin{figure}[h]
\centering
\begin{tikzpicture}[font=\footnotesize, >=latex]
  \foreach \x/\n in {0/Cliente, 3/{Nodo1 (Líder)}, 6.5/{Nodo2 (Seguidor)}, 10/{Nodo3 (Seguidor)}} {
    \node[draw, fill=blue!10, minimum width=2.2cm] at (\x,0) (\n) {\n};
    \draw[dashed] (\x,-0.4) -- (\x,-7);
  }
  \draw[->] (0,-1) -- node[above]{escritura} (3,-1);
  \draw[->] (3,-1.8) -- node[above]{append log} (6.5,-1.8);
  \draw[->] (3,-2.3) -- node[above]{append log} (10,-2.3);
  \draw[<-] (3,-3) -- node[above]{ack} (6.5,-3);
  \draw[<-] (3,-3.5) -- node[above]{ack} (10,-3.5);
  \node[red] at (6.5,-4) {\Large